\IfFileExists{stacks-project.cls}{%
\documentclass{stacks-project}
}{%
\documentclass{amsart}
}

% For dealing with references we use the comment environment
\usepackage{verbatim}
\newenvironment{reference}{\comment}{\endcomment}
%\newenvironment{reference}{}{}
\newenvironment{slogan}{\comment}{\endcomment}
\newenvironment{history}{\comment}{\endcomment}

% For commutative diagrams we use Xy-pic
\usepackage[all]{xy}

% We use 2cell for 2-commutative diagrams.
\xyoption{2cell}
\UseAllTwocells

% We use multicol for the list of chapters between chapters
\usepackage{multicol}

% This is generally recommended for better output
\usepackage{lmodern}
\usepackage[T1]{fontenc}

% For cross-file-references
\usepackage{xr-hyper}

% Package for hypertext links:
\usepackage{hyperref}

% For any local file, say "hello.tex" you want to link to please
% use \externaldocument[hello-]{hello}
\externaldocument[introduction-]{introduction}
\externaldocument[conventions-]{conventions}
\externaldocument[sets-]{sets}
\externaldocument[categories-]{categories}
\externaldocument[topology-]{topology}
\externaldocument[sheaves-]{sheaves}
\externaldocument[sites-]{sites}
\externaldocument[stacks-]{stacks}
\externaldocument[fields-]{fields}
\externaldocument[algebra-]{algebra}
\externaldocument[brauer-]{brauer}
\externaldocument[homology-]{homology}
\externaldocument[derived-]{derived}
\externaldocument[simplicial-]{simplicial}
\externaldocument[more-algebra-]{more-algebra}
\externaldocument[smoothing-]{smoothing}
\externaldocument[modules-]{modules}
\externaldocument[sites-modules-]{sites-modules}
\externaldocument[injectives-]{injectives}
\externaldocument[cohomology-]{cohomology}
\externaldocument[sites-cohomology-]{sites-cohomology}
\externaldocument[dga-]{dga}
\externaldocument[dpa-]{dpa}
\externaldocument[sdga-]{sdga}
\externaldocument[hypercovering-]{hypercovering}
\externaldocument[schemes-]{schemes}
\externaldocument[constructions-]{constructions}
\externaldocument[properties-]{properties}
\externaldocument[morphisms-]{morphisms}
\externaldocument[coherent-]{coherent}
\externaldocument[divisors-]{divisors}
\externaldocument[limits-]{limits}
\externaldocument[varieties-]{varieties}
\externaldocument[topologies-]{topologies}
\externaldocument[descent-]{descent}
\externaldocument[perfect-]{perfect}
\externaldocument[more-morphisms-]{more-morphisms}
\externaldocument[flat-]{flat}
\externaldocument[groupoids-]{groupoids}
\externaldocument[more-groupoids-]{more-groupoids}
\externaldocument[etale-]{etale}
\externaldocument[chow-]{chow}
\externaldocument[intersection-]{intersection}
\externaldocument[pic-]{pic}
\externaldocument[weil-]{weil}
\externaldocument[adequate-]{adequate}
\externaldocument[dualizing-]{dualizing}
\externaldocument[duality-]{duality}
\externaldocument[discriminant-]{discriminant}
\externaldocument[derham-]{derham}
\externaldocument[local-cohomology-]{local-cohomology}
\externaldocument[algebraization-]{algebraization}
\externaldocument[curves-]{curves}
\externaldocument[resolve-]{resolve}
\externaldocument[models-]{models}
\externaldocument[functors-]{functors}
\externaldocument[equiv-]{equiv}
\externaldocument[pione-]{pione}
\externaldocument[etale-cohomology-]{etale-cohomology}
\externaldocument[proetale-]{proetale}
\externaldocument[relative-cycles-]{relative-cycles}
\externaldocument[more-etale-]{more-etale}
\externaldocument[trace-]{trace}
\externaldocument[crystalline-]{crystalline}
\externaldocument[spaces-]{spaces}
\externaldocument[spaces-properties-]{spaces-properties}
\externaldocument[spaces-morphisms-]{spaces-morphisms}
\externaldocument[decent-spaces-]{decent-spaces}
\externaldocument[spaces-cohomology-]{spaces-cohomology}
\externaldocument[spaces-limits-]{spaces-limits}
\externaldocument[spaces-divisors-]{spaces-divisors}
\externaldocument[spaces-over-fields-]{spaces-over-fields}
\externaldocument[spaces-topologies-]{spaces-topologies}
\externaldocument[spaces-descent-]{spaces-descent}
\externaldocument[spaces-perfect-]{spaces-perfect}
\externaldocument[spaces-more-morphisms-]{spaces-more-morphisms}
\externaldocument[spaces-flat-]{spaces-flat}
\externaldocument[spaces-groupoids-]{spaces-groupoids}
\externaldocument[spaces-more-groupoids-]{spaces-more-groupoids}
\externaldocument[bootstrap-]{bootstrap}
\externaldocument[spaces-pushouts-]{spaces-pushouts}
\externaldocument[spaces-chow-]{spaces-chow}
\externaldocument[groupoids-quotients-]{groupoids-quotients}
\externaldocument[spaces-more-cohomology-]{spaces-more-cohomology}
\externaldocument[spaces-simplicial-]{spaces-simplicial}
\externaldocument[spaces-duality-]{spaces-duality}
\externaldocument[formal-spaces-]{formal-spaces}
\externaldocument[restricted-]{restricted}
\externaldocument[spaces-resolve-]{spaces-resolve}
\externaldocument[formal-defos-]{formal-defos}
\externaldocument[defos-]{defos}
\externaldocument[cotangent-]{cotangent}
\externaldocument[examples-defos-]{examples-defos}
\externaldocument[algebraic-]{algebraic}
\externaldocument[examples-stacks-]{examples-stacks}
\externaldocument[stacks-sheaves-]{stacks-sheaves}
\externaldocument[criteria-]{criteria}
\externaldocument[artin-]{artin}
\externaldocument[quot-]{quot}
\externaldocument[stacks-properties-]{stacks-properties}
\externaldocument[stacks-morphisms-]{stacks-morphisms}
\externaldocument[stacks-limits-]{stacks-limits}
\externaldocument[stacks-cohomology-]{stacks-cohomology}
\externaldocument[stacks-perfect-]{stacks-perfect}
\externaldocument[stacks-introduction-]{stacks-introduction}
\externaldocument[stacks-more-morphisms-]{stacks-more-morphisms}
\externaldocument[stacks-geometry-]{stacks-geometry}
\externaldocument[moduli-]{moduli}
\externaldocument[moduli-curves-]{moduli-curves}
\externaldocument[examples-]{examples}
\externaldocument[exercises-]{exercises}
\externaldocument[guide-]{guide}
\externaldocument[desirables-]{desirables}
\externaldocument[coding-]{coding}
\externaldocument[obsolete-]{obsolete}
\externaldocument[fdl-]{fdl}
\externaldocument[index-]{index}

% Theorem environments.
%
\theoremstyle{plain}
\newtheorem{theorem}[subsection]{Theorem}
\newtheorem{proposition}[subsection]{Proposition}
\newtheorem{lemma}[subsection]{Lemma}

\theoremstyle{definition}
\newtheorem{definition}[subsection]{Definition}
\newtheorem{example}[subsection]{Example}
\newtheorem{exercise}[subsection]{Exercise}
\newtheorem{situation}[subsection]{Situation}

\theoremstyle{remark}
\newtheorem{remark}[subsection]{Remark}
\newtheorem{remarks}[subsection]{Remarks}

\numberwithin{equation}{subsection}

% Macros
%
\def\lim{\mathop{\mathrm{lim}}\nolimits}
\def\colim{\mathop{\mathrm{colim}}\nolimits}
\def\Spec{\mathop{\mathrm{Spec}}}
\def\Hom{\mathop{\mathrm{Hom}}\nolimits}
\def\Ext{\mathop{\mathrm{Ext}}\nolimits}
\def\SheafHom{\mathop{\mathcal{H}\!\mathit{om}}\nolimits}
\def\SheafExt{\mathop{\mathcal{E}\!\mathit{xt}}\nolimits}
\def\Sch{\mathit{Sch}}
\def\Mor{\mathop{\mathrm{Mor}}\nolimits}
\def\Ob{\mathop{\mathrm{Ob}}\nolimits}
\def\Sh{\mathop{\mathit{Sh}}\nolimits}
\def\NL{\mathop{N\!L}\nolimits}
\def\CH{\mathop{\mathrm{CH}}\nolimits}
\def\proetale{{pro\text{-}\acute{e}tale}}
\def\etale{{\acute{e}tale}}
\def\QCoh{\mathit{QCoh}}
\def\Ker{\mathop{\mathrm{Ker}}}
\def\Im{\mathop{\mathrm{Im}}}
\def\Coker{\mathop{\mathrm{Coker}}}
\def\Coim{\mathop{\mathrm{Coim}}}

% Boxtimes
%
\DeclareMathSymbol{\boxtimes}{\mathbin}{AMSa}{"02}

%
% Macros for moduli stacks/spaces
%
\def\QCohstack{\mathcal{QC}\!\mathit{oh}}
\def\Cohstack{\mathcal{C}\!\mathit{oh}}
\def\Spacesstack{\mathcal{S}\!\mathit{paces}}
\def\Quotfunctor{\mathrm{Quot}}
\def\Hilbfunctor{\mathrm{Hilb}}
\def\Curvesstack{\mathcal{C}\!\mathit{urves}}
\def\Polarizedstack{\mathcal{P}\!\mathit{olarized}}
\def\Complexesstack{\mathcal{C}\!\mathit{omplexes}}
% \Pic is the operator that assigns to X its picard group, usage \Pic(X)
% \Picardstack_{X/B} denotes the Picard stack of X over B
% \Picardfunctor_{X/B} denotes the Picard functor of X over B
\def\Pic{\mathop{\mathrm{Pic}}\nolimits}
\def\Picardstack{\mathcal{P}\!\mathit{ic}}
\def\Picardfunctor{\mathrm{Pic}}
\def\Deformationcategory{\mathcal{D}\!\mathit{ef}}
\usepackage{cleveref}
\crefname{originaldefinition}{original definition}{original definitions}
\Crefname{originaldefinition}{Original definition}{Original definitions}
\makeatletter
\newlabel{definition-divided-powers-dga}{{09PJ}{source}{}{sourcepointer-1}{}}
\newlabel{definition-divided-powers-dga@cref}{{[originaldefinition][0][]09PJ}{[1][0][]source}{}{}{}}
\makeatother
\crefname{originalreference}{original reference}{original references}
\Crefname{originalreference}{Original reference}{Original references}
\makeatletter
\newlabel{example-adjoining-even}{{09PI}{source}{}{sourcepointer-2}{}}
\newlabel{example-adjoining-even@cref}{{[originalreference][0][]09PI}{[1][0][]source}{}{}{}}
\makeatother
\crefname{originalreference}{original reference}{original references}
\Crefname{originalreference}{Original reference}{Original references}
\makeatletter
\newlabel{example-adjoining-odd}{{09PH}{source}{}{sourcepointer-3}{}}
\newlabel{example-adjoining-odd@cref}{{[originalreference][0][]09PH}{[1][0][]source}{}{}{}}
\makeatother
\crefname{originallemma}{original lemma}{original lemmas}
\Crefname{originallemma}{Original lemma}{Original lemmas}
\makeatletter
\newlabel{lemma-tate-resolution}{{09PN}{source}{}{sourcepointer-4}{}}
\newlabel{lemma-tate-resolution@cref}{{[originallemma][0][]09PN}{[1][0][]source}{}{}{}}
\makeatother
\crefname{originalreference}{original reference}{original references}
\Crefname{originalreference}{Original reference}{Original references}
\makeatletter
\newlabel{more-algebra-lemma-sequence-Koszul-complexes}{{0921}{source}{}{sourcepointer-5}{}}
\newlabel{more-algebra-lemma-sequence-Koszul-complexes@cref}{{[originalreference][0][]0921}{[1][0][]source}{}{}{}}
\makeatother
\title{Lifting maps between the Tate resolution and Koszul towers}
\author{The Stacks Project Authors}
\date{stacks-project:dpa.tex:local-snapshot-20261012; GNU FDL 1.2 or later}
\begin{document}
\maketitle
\noindent\textit{Editorial note.} 'The map phi\_2 exists and is compatible with the differential and the divided powers by the universal [property] of the divided power polynomial algebra' -- the universal-property invocation is in-chapter template (lemma-divided-power-polynomial-algebra); no other omission markers.\par\medskip
\noindent\textit{Allowed standard prerequisites (editorial).}\begin{itemize}
\item Tate resolution construction. lemma-tate-resolution (same chapter, named). Supplies the inductive skeleton this proof reuses ('exactly as in the statement and proof of Lemma lemma-tate-resolution').
\item Koszul transition selection. more-algebra-lemma-sequence-Koszul-complexes (cross-chapter, named). Chooses the tower n = n\_0 < n\_1 < ... < n\_r = N with zero maps on positive homology, the spine of the construction.
\end{itemize}\medskip
\ifdefined\hypertarget\hypertarget{sourcepointer-1}{}\fi
\noindent\textit{Source reference.} Original definition 09PJ: sources/stacks-dpa/pinned-original-source.tex, source lines 1300--1300. Navigation pointer only: the counted proof cites this statement by \textbackslash{}ref as a standing input; the deduction the counted proof has to make is carried out inside the counted span, and no part of that deduction is delegated to this pointer.\par\medskip
\ifdefined\hypertarget\hypertarget{sourcepointer-2}{}\fi
\noindent\textit{Source reference.} Original reference 09PI: sources/stacks-dpa/pinned-original-source.tex, source lines 1241--1241. Navigation pointer only: the counted proof cites this statement by \textbackslash{}ref as a standing input; the deduction the counted proof has to make is carried out inside the counted span, and no part of that deduction is delegated to this pointer.\par\medskip
\ifdefined\hypertarget\hypertarget{sourcepointer-3}{}\fi
\noindent\textit{Source reference.} Original reference 09PH: sources/stacks-dpa/pinned-original-source.tex, source lines 1221--1221. Navigation pointer only: the counted proof cites this statement by \textbackslash{}ref as a standing input; the deduction the counted proof has to make is carried out inside the counted span, and no part of that deduction is delegated to this pointer.\par\medskip
\ifdefined\hypertarget\hypertarget{sourcepointer-4}{}\fi
\noindent\textit{Source reference.} Original lemma 09PN: sources/stacks-dpa/pinned-original-source.tex, source lines 1424--1424. Navigation pointer only: the counted proof cites this statement by \textbackslash{}ref as a standing input; the deduction the counted proof has to make is carried out inside the counted span, and no part of that deduction is delegated to this pointer.\par\medskip
\ifdefined\hypertarget\hypertarget{sourcepointer-5}{}\fi
\noindent\textit{Source reference.} Original reference 0921: focused-seeds/geometry/S9-dpa-lift-tate-to-koszul/ref-tags.txt, source lines 5--5. Navigation pointer only: the counted proof cites this statement by \textbackslash{}ref as a standing input; the deduction the counted proof has to make is carried out inside the counted span, and no part of that deduction is delegated to this pointer.\par\medskip
\begin{lemma}
\label{lemma-lift-tate-to-koszul}
In the situation above, if $R$ is Noetherian, then
for every $n$ there exists an $N \geq n$ and maps
$$
K_N \to A \to R/(f_1^N, \ldots, f_r^N)\quad\text{and}\quad A \to K_n
$$
with the following properties
\begin{enumerate}
\item $(A, \text{d}, \gamma)$ is as in
Definition \ref{definition-divided-powers-dga},
\item $A \to R/(f_1^N, \ldots, f_r^N)$ is a quasi-isomorphism,
\item the composition $K_N \to A \to R/(f_1^N, \ldots, f_r^N)$
is the canonical map,
\item the composition $K_N \to A \to K_n$ is the transition map,
\item $A_0 = R \to R/(f_1^N, \ldots, f_r^N)$ is the canonical
surjection,
\item $A$ is a graded divided power polynomial algebra over $R$
with finitely many generators in each degree, and
\item $A \to K_n$ is a homomorphism of differential graded $R$-algebras
compatible with divided powers which induces the canonical map
$R/(f_1^N, \ldots, f_r^N) \to R/(f_1^n, \ldots, f_r^n)$ on
homology in degree $0$.
\end{enumerate}
Condition (4) means that $A$ is constructed out of $A_0$ by
successively adjoining a finite set of variables $T$ in each degree
$> 0$ as in Example \ref{example-adjoining-odd} or \ref{example-adjoining-even}.
\end{lemma}

\begin{proof}
Fix $n$. If $r = 0$, then we can just pick $A = R$. Assume $r > 0$. By
More on Algebra, Lemma \ref{more-algebra-lemma-sequence-Koszul-complexes}
(translated into the language of chain complexes) we can choose
$$
n_{r} > n_{r - 1} > \ldots > n_1 > n_0 = n
$$
such that the transition maps $K_{n_{i + 1}} \to K_{n_i}$ on Koszul
algebras (see above) induce the zero map on homology in degrees $> 0$.
We will prove the lemma with $N = n_r$.

\medskip\noindent
We will construct $A$ exactly as in the statement and proof of
Lemma \ref{lemma-tate-resolution}. Thus we will have
$$
A = \colim A(m),\quad\text{and}\quad
A(0) \to A(1) \to A(2) \to \ldots \to R/(f_1^N, \ldots, f_r^N)
$$
This will immediately give us properties (1), (2), (5), and (6).
To finish the proof we will construct the $R$-algebra maps
$K_N \to A \to K_n$. To do this we will construct
\begin{enumerate}
\item an isomorphism $A(1) \to K_N = K_{n_r}$,
\item a map $A(2) \to K_{n_{r - 1}}$,
\item $\ldots$
\item a map $A(r) \to K_{n_1}$,
\item a map $A(r + 1) \to K_{n_0} = K_n$, and
\item a map $A \to K_n$.
\end{enumerate}
In each of these steps the map constructed will be between
differential graded algebras compatibly endowed with divided powers
and each of the maps will be compatible with the
previous one via the transition maps between the Koszul algebras
and each of the maps will induce the obvious canonical map
on homology in degree $0$.

\medskip\noindent
Recall that $A(0) = R$. For $m = 1$, the proof of
Lemma \ref{lemma-tate-resolution}
chooses $A(1) = R\langle T_1, \ldots, T_r\rangle$ with
$T_i$ of degree $1$ and with $\text{d}(T_i) = f_i^N$.
Namely, the $f_i^N$ are generators of the kernel of
$A(0) \to R/(f_1^N, \ldots, f_r^N)$.
Thus for $A(1) \to K_N = K_{n_r}$ we use the map
$$
\varphi_1 : A(1) \longrightarrow K_{n_r},\quad T_i \longmapsto \xi_i
$$
which is an isomorphism.

\medskip\noindent
For $m = 2$, the construction in the proof of Lemma \ref{lemma-tate-resolution}
chooses generators $e_1, \ldots, e_t \in \Ker(\text{d} : A(1)_1 \to A(1)_0)$.
The construction proceeds by taking
$A(2) = A(1)\langle T_1, \ldots, T_t\rangle$
as a divided power polynomial algebra with $T_i$ of degree $2$
and with $\text{d}(T_i) = e_i$.
Since $\varphi_1(e_i)$ is a cocycle in $K_{n_r}$
we see that its image in $K_{n_{r - 1}}$ is a coboundary by
our choice of $n_r$ and $n_{r - 1}$ above.
Hence we can construct the following commutative diagram
$$
\xymatrix{
A(1) \ar[d] \ar[r]_{\varphi_1} & K_{n_r} \ar[d] \\
A(2) \ar[r]^{\varphi_2} & K_{n_{r - 1}}
}
$$
by sending $T_i$ to an element in degree $2$ whose boundary is the
image of $\varphi_1(e_i)$. The map $\varphi_2$ exists and is compatible
with the differential and the divided powers by the universal
of the divided power polynomial algebra.

\medskip\noindent
The algebra $A(m)$ and the map $\varphi_m : A(m) \to K_{n_{r + 1 - m}}$
are constructed in exactly the same manner for $m = 2, \ldots, r$.

\medskip\noindent
Given the map $A(r) \to K_{n_1}$ we see that the composition
$H_r(A(r)) \to H_r(K_{n_1}) \to H_r(K_{n_0}) \subset (K_{n_0})_r$
is zero, hence we can extend this to $A(r + 1) \to K_{n_0} = K_n$
by sending the new polynomial generators of $A(r + 1)$ to zero.

\medskip\noindent
Having constructed $A(r + 1) \to K_{n_0} = K_n$ we can simply
extend to $A(r + 2), A(r + 3), \ldots$ in the only possible way
by sending the new polynomial generators to zero.
This finishes the proof.
\end{proof}

\end{document}
